% !TeX root = ../main.tex
\chapter{THỦ TỤC – HÀM - TRIGGER}\label{chap:routines}
\section{Thủ tục INSERT / UPDATE / DELETE dữ liệu vào 1 bảng dữ liệu}
\subsection{Nguyên tắc xử lý và bảng dữ liệu chính}
Bảng dữ liệu chính liên quan: \code{KHU_VUC}. Bảng dữ liệu phụ liên quan: \code{MAY_TRAM}, \code{PHAN_CONG_TRUC}. Thủ tục kiểm tra đầy đủ điều kiện ràng buộc nghiệp vụ và xử lý ngoại lệ có mã lỗi rõ ràng. Mỗi thủ tục trả về mã khu vực, tên thao tác và số dòng bị tác động.
\subsection{Thủ tục INSERT: Thêm khu vực (sp\_khu\_vuc\_insert)}
Mô tả: Thêm khu vực mới trong màn hình quản trị mặt bằng quán. Kiểm tra mã không rỗng, tên không trùng. Input: \code{p_ma_khu_vuc}, \code{p_ten_khu_vuc}, \code{p_mo_ta_vi_tri}. Output: \code{ma_khu_vuc}, \code{thao_tac}, \code{so_dong_tac_dong}.
\flowfigure{diagrams/crud_insert.tex}{Lưu đồ minh họa thao tác xử lý thủ tục INSERT}{fig:flow-crud_insert}
\sqlinput{sql/insert.sql}{Câu lệnh tạo thủ tục INSERT}
\sqlinput{sql/demo_img_09_crud_insert.sql}{Câu lệnh kiểm chứng trên MySQL}
\snapshot{img_09_crud_insert}{Kết quả từ DBMS minh họa kiểm tra thủ tục INSERT (trường hợp báo lỗi và trường hợp đúng)}{fig:img_09_crud_insert}
\subsection{Thủ tục UPDATE: Cập nhật khu vực (sp\_khu\_vuc\_update)}
Mô tả: Cập nhật thông tin tên và vị trí khu vực. Kiểm tra khu vực tồn tại, tên mới không trùng khu vực khác. Input: \code{p_ma_khu_vuc}, \code{p_ten_moi}, \code{p_mo_ta_moi}. Output: \code{ma_khu_vuc}, \code{thao_tac}, \code{so_dong_tac_dong}.
\flowfigure{diagrams/crud_update.tex}{Lưu đồ minh họa thao tác xử lý thủ tục UPDATE}{fig:flow-crud_update}
\sqlinput{sql/update.sql}{Câu lệnh tạo thủ tục UPDATE}
\sqlinput{sql/demo_img_10_crud_update.sql}{Câu lệnh kiểm chứng trên MySQL}
\snapshot{img_10_crud_update}{Kết quả từ DBMS minh họa kiểm tra thủ tục UPDATE (trường hợp báo lỗi và trường hợp đúng)}{fig:img_10_crud_update}
\subsection{Thủ tục DELETE: Xóa khu vực (sp\_khu\_vuc\_delete)}
Mô tả: Xóa khu vực không còn hoạt động. Kiểm tra khu vực tồn tại và không còn máy trạm hoặc lịch trực tham chiếu. Input: \code{p_ma_khu_vuc}. Output: \code{ma_khu_vuc}, \code{thao_tac}, \code{so_dong_tac_dong}.
\flowfigure{diagrams/crud_delete.tex}{Lưu đồ minh họa thao tác xử lý thủ tục DELETE}{fig:flow-crud_delete}
\sqlinput{sql/delete.sql}{Câu lệnh tạo thủ tục DELETE}
\sqlinput{sql/demo_img_11_crud_delete.sql}{Câu lệnh kiểm chứng trên MySQL}
\snapshot{img_11_crud_delete}{Kết quả từ DBMS minh họa kiểm tra thủ tục DELETE (trường hợp báo lỗi và trường hợp đúng)}{fig:img_11_crud_delete}
\snapshot{routines_created}{Kết quả màn hình hiển thị từ DBMS minh họa cho việc tạo các thủ tục CRUD}{fig:routines_created}
\section{TRIGGER}
\subsection{Trigger 1: Bảo vệ biên lai đã thanh toán (T-P12-01)}
Bảng dữ liệu chính liên quan: \code{BIEN_LAI_TIEN_MAY}. Mô tả: Trigger kiểm tra ràng buộc nghiệp vụ kế toán — không cho phép sửa đổi hoặc xóa biên lai đã ở trạng thái thanh toán (\code{DA_THANH_TOAN}). Được kích hoạt BEFORE UPDATE và BEFORE DELETE. Vẽ flowchart minh họa xử lý.
\flowfigure{diagrams/receipt.tex}{Lưu đồ minh họa xử lý Trigger 1 bảo vệ biên lai đã thanh toán}{fig:flow-receipt}
\sqlinput{sql/trigger_receipt.sql}{Câu lệnh tạo Trigger 1 trên BIEN\_LAI\_TIEN\_MAY}
\sqlinput{sql/demo_img_12_trigger_t1.sql}{Câu lệnh kiểm chứng trên MySQL}
\snapshot{img_12_trigger_t1}{Kết quả DBMS minh họa trường hợp báo lỗi: Chặn sửa và xóa biên lai đã thanh toán}{fig:img_12_trigger_t1}
\sqlinput{sql/demo_trigger_t1_valid.sql}{Câu lệnh kiểm chứng trên MySQL}
\snapshot{trigger_t1_valid}{Kết quả DBMS minh họa trường hợp đúng: Cho phép sửa và xóa biên lai chưa thanh toán}{fig:trigger_t1_valid}
\subsection{Trigger 2: Đồng bộ tổng tiền đơn từ chi tiết (T-P12-02)}
Bảng dữ liệu chính liên quan: \code{CHI_TIET_GOI_DICH_VU}. Bảng dữ liệu phụ liên quan: \code{DON_GOI_DICH_VU}. Mô tả: Nhóm trigger thao tác ảnh hưởng tới hơn 1 bảng — tự động cập nhật tổng tiền đầu đơn (\code{tong_tien}) khi thêm, sửa số lượng/giá, hoặc xóa dòng chi tiết món. Kích hoạt AFTER INSERT, AFTER UPDATE, AFTER DELETE.
\flowfigure{diagrams/order_total.tex}{Lưu đồ minh họa xử lý Trigger 2 đồng bộ tổng giá trị đơn hàng}{fig:flow-order_total}
\sqlinput{sql/trigger_total.sql}{Câu lệnh tạo nhóm Trigger 2 trên CHI\_TIET\_GOI\_DICH\_VU}
\sqlinput{sql/demo_trigger_t2_before.sql}{Câu lệnh kiểm chứng trên MySQL}
\snapshot{trigger_t2_before}{Dữ liệu đơn hàng ban đầu có tổng tiền bằng 0}{fig:trigger_t2_before}
\sqlinput{sql/demo_img_13_trigger_t2.sql}{Câu lệnh kiểm chứng trên MySQL}
\snapshot{img_13_trigger_t2}{Kết quả DBMS minh họa trường hợp đúng: Tổng tiền tự động đồng bộ khi thêm/sửa/xóa dòng món}{fig:img_13_trigger_t2}
\sqlinput{sql/demo_trigger_t2_errors.sql}{Câu lệnh kiểm chứng trên MySQL}
\snapshot{trigger_t2_errors}{Kết quả DBMS minh họa trường hợp báo lỗi: Chặn thao tác trên đơn đã đóng hoặc đã hủy}{fig:trigger_t2_errors}
\snapshot{triggers_created}{Các trigger trong hai yêu cầu đã được tạo}{fig:triggers_created}
\section{THỦ TỤC}
\subsection{Thủ tục 1: Báo cáo doanh thu theo ngày (sp\_bao\_cao\_doanh\_thu\_ngay)}
Mô tả thủ tục: Tổng hợp dữ liệu từ nhiều bảng (\code{BIEN_LAI_TIEN_MAY}, \code{DON_GOI_DICH_VU}). Sử dụng câu lệnh điều kiện IF kiểm tra ngày hợp lệ và vòng lặp WHILE duyệt qua từng ngày trong khoảng thời gian để tính tổng thu máy, thu dịch vụ và hoàn tiền. Output: Mảng JSON doanh thu chi tiết từng ngày.
\flowfigure{diagrams/revenue.tex}{Lưu đồ minh họa Thủ tục 1 duyệt vòng lặp WHILE tính doanh thu theo ngày}{fig:flow-revenue}
\sqlinput{sql/report_revenue.sql}{Câu lệnh tạo Thủ tục 1: sp\_bao\_cao\_doanh\_thu\_ngay}
\sqlinput{sql/demo_img_14_sp_doanh_thu.sql}{Câu lệnh kiểm chứng trên MySQL}
\snapshot{img_14_sp_doanh_thu}{Kết quả DBMS minh họa kiểm tra Thủ tục 1 (ca báo lỗi ngày sai và ca đúng trả mảng JSON)}{fig:img_14_sp_doanh_thu}
\subsection{Thủ tục 2: Báo cáo tồn và phân loại dịch vụ (sp\_bao\_cao\_ton\_dich\_vu)}
Mô tả thủ tục: Đọc dữ liệu từ bảng \code{DICH_VU}, \code{NHOM_DICH_VU}, \code{DON_VI_TINH}. Sử dụng CURSOR duyệt danh sách món kết hợp IF/ELSE phân loại trạng thái tồn kho (\code{HetHang}, \code{SapHet}, \code{DuHang}) theo ngưỡng chỉ định. Output: Mảng JSON danh sách dịch vụ và trạng thái phân loại.
\flowfigure{diagrams/stock.tex}{Lưu đồ minh họa Thủ tục 2 sử dụng CURSOR và IF phân loại tồn kho}{fig:flow-stock}
\sqlinput{sql/report_stock.sql}{Câu lệnh tạo Thủ tục 2: sp\_bao\_cao\_ton\_dich\_vu}
\sqlinput{sql/demo_img_15_sp_ton_kho.sql}{Câu lệnh kiểm chứng trên MySQL}
\snapshot{img_15_sp_ton_kho}{Kết quả DBMS minh họa kiểm tra Thủ tục 2 (ca báo lỗi ngưỡng âm và ca đúng trả dữ liệu)}{fig:img_15_sp_ton_kho}
\section{HÀM}
\subsection{Hàm 1: Tổng chi ví ròng (fn\_tong\_chi\_vi\_rong)}
Mô tả hàm: Đọc dữ liệu từ bảng \code{GIAO_DICH_VI} và \code{NGUOI_CHOI}. Sử dụng CURSOR duyệt qua lịch sử giao dịch kết hợp câu lệnh IF kiểm tra loại giao dịch (\code{TruTienMay}, \code{TruTienDichVu}, \code{HoanTien}) để tính tổng số tiền chi ròng của hội viên trong một khoảng thời gian. Input: \code{p_ma_nguoi_choi}, \code{p_tu_gio}, \code{p_den_gio}. Output: DECIMAL(18,0).
\flowfigure{diagrams/wallet.tex}{Lưu đồ minh họa Hàm 1 dùng CURSOR và IF tính tổng chi ròng}{fig:flow-wallet}
\sqlinput{sql/fn_wallet.sql}{Câu lệnh tạo Hàm 1: fn\_tong\_chi\_vi\_rong}
\sqlinput{sql/demo_img_16_fn_chi_vi.sql}{Câu lệnh kiểm chứng trên MySQL}
\snapshot{img_16_fn_chi_vi}{Kết quả DBMS minh họa kiểm tra Hàm 1 (ca đúng trả tổng tiền và ca sai mã trả NULL)}{fig:img_16_fn_chi_vi}
\subsection{Hàm 2: Tổng tiền từ dòng chi tiết đơn (fn\_tong\_tien\_don)}
Mô tả hàm: Đọc dữ liệu từ bảng \code{CHI_TIET_GOI_DICH_VU}. Sử dụng CURSOR duyệt từng dòng chi tiết của đơn hàng kết hợp kiểm tra IF điều kiện số lượng và đơn giá hợp lệ, sau đó cộng dồn thành tiền của toàn bộ đơn. Input: \code{p_ma_don_goi}. Output: DECIMAL(18,0).
\flowfigure{diagrams/order.tex}{Lưu đồ minh họa Hàm 2 dùng CURSOR tính tổng tiền đơn hàng}{fig:flow-order}
\sqlinput{sql/fn_order.sql}{Câu lệnh tạo Hàm 2: fn\_tong\_tien\_don}
\sqlinput{sql/demo_img_17_fn_tien_don.sql}{Câu lệnh kiểm chứng trên MySQL}
\snapshot{img_17_fn_tien_don}{Kết quả DBMS minh họa kiểm tra Hàm 2 (ca đúng trả tổng tiền và ca sai mã trả NULL)}{fig:img_17_fn_tien_don}
